\documentclass{article}
\usepackage{amsmath}
\usepackage{amsfonts}

\title{Mathematical Formulations of A.L.I.V.E. NEXUS}
\author{Devanik}
\date{April 2026}

\begin{document}

\maketitle

\section{Introduction}
This document formalizes the mathematical underpinnings of the A.L.I.V.E. NEXUS project.

\section{Reinforcement Learning Engine}
\subsection{Dueling DQN}
The Q-value function is decomposed into a value stream $V(s)$ and an advantage stream $A(s, a)$:
\begin{equation}
Q(s, a; \theta, \alpha, \beta) = V(s; \theta, \beta) + \left( A(s, a; \theta, \alpha) - \frac{1}{|\mathcal{A}|} \sum_{a'} A(s, a'; \theta, \alpha) \right)
\end{equation}

\subsection{N-Step Returns}
The $N$-step return is given by:
\begin{equation}
G_t^{(n)} = \sum_{k=0}^{n-1} \gamma^k r_{t+k} + \gamma^n \max_{a'} Q_{\text{target}}(s_{t+n}, a')
\end{equation}

\subsection{Prioritized Experience Replay}
Sampling probability for transition $i$:
\begin{equation}
P(i) = \frac{p_i^{\alpha}}{\sum_j p_j^{\alpha}}
\end{equation}
With importance-sampling weights:
\begin{equation}
w_i = \left( \frac{1}{N \cdot P(i)} \right)^{\beta}
\end{equation}

\end{document}
